\documentclass[11pt]{article}

\usepackage[margin=1in]{geometry}
\usepackage{hyperref}
\usepackage{url}

\title{Capstone --- Lost and Found\\\large Citations}
\author{Jason Shan \and Mustafa Kurt \and Raymond Huang}
\date{\today}

\begin{document}
\maketitle

\section*{External contributions}

This document lists external contributions that influenced the
\emph{Capstone --- Lost and Found} project, including external APIs,
AI tools used during development, and assets sourced from the web.
Packages declared in dependency files (e.g.\ \texttt{package.json}) are
not individually cited here, in line with the project requirements.

\subsection*{External APIs and SDKs}

The application integrates the Anthropic Messages
API~\cite{anthropic_api} to power the in-app Hawk~AI chatbot,
Supabase~\cite{supabase} for database and file storage, and
Socket.IO~\cite{socketio} for real-time student--admin messaging.
The user interface uses the \emph{Work Sans} and \emph{Oswald}
typefaces from Google Fonts~\cite{googlefonts}.

\subsection*{Hosting platforms}

The React/Vite frontend is hosted on Vercel~\cite{vercel}, and the
Node/Express backend (REST API and Socket.IO server) is hosted on
Render~\cite{render}.

\subsection*{AI tools used during development}

Per the project requirements, the Claude model that powers the in-app
chatbot~\cite{ai_claude_chatbot} is cited separately from the Claude
instance used as a development assistant~\cite{ai_claude_assistant}.
In addition, the team used ChatGPT~\cite{ai_chatgpt} and the
Cursor~\cite{ai_cursor} AI-assisted code editor during development for
questions, debugging help, and code suggestions.

\subsection*{Design inspiration}

The visual layout, color palette, and overall page structure of the
Lost-and-Found web application were modeled after the official Hunter
College website~\cite{hunter_college_site} so that the application
would feel consistent with the school's existing web presence. No
code was copied; only the visual design language was referenced.

\subsection*{Assets}

Campus photographs and mock lost-item images shipped under
\texttt{client/src/assets/} were sourced via Google Image Search to
simulate realistic lost-and-found scenarios for demonstration
purposes~\cite{assets_images}.

\bibliographystyle{plain}
\bibliography{references}

\end{document}
